\chapter{Tensor-Network States}
\label{ch:lec15}

\begin{objectives}
\begin{itemize}[leftmargin=1.4em,nosep]
  \item Read and manipulate tensor-network diagrams.
  \item Derive matrix product states (MPS) from sequential Schmidt decompositions.
  \item Put an MPS into canonical form and bound its contraction cost.
  \item Relate bond dimension to entanglement and the area law.
\end{itemize}
\end{objectives}

\todo{Write the introductory paragraph orienting the lecture.}

\section{Tensors and diagrammatic notation}
\label{sec:lec15-1}
\todo{Draft this section.}

\section{Schmidt decomposition and the SVD}
\label{sec:lec15-2}
\todo{Draft this section.}

\section{Matrix product states}
\label{sec:lec15-3}
\todo{Draft this section.}

\section{Canonical forms and gauge freedom}
\label{sec:lec15-4}
\todo{Draft this section.}

\section{Contraction cost, bond dimension, and the area law}
\label{sec:lec15-5}
\todo{Draft this section.}

\exercisesection
\begin{exercise}
\todo{Add exercise.}
\end{exercise}
\begin{exercise}
\todo{Add exercise.}
\end{exercise}
\begin{exercise}
\todo{Add exercise.}
\end{exercise}

\readingsection
\begin{itemize}[leftmargin=1.4em]
  \item U.~Schollw\\\"ock, \emph{The density-matrix renormalization group in the age of matrix product states}, Ann.~Phys.~326 (2011).
  \item R.~Or\\'us, \emph{A practical introduction to tensor networks}, Ann.~Phys.~349 (2014).
\end{itemize}
